\documentclass[11pt,a4paper]{article}
\usepackage[utf8]{inputenc}
\usepackage[T1]{fontenc}
\usepackage{lmodern}
\usepackage[french]{babel}
\usepackage{amsmath,amssymb,mathtools}
\usepackage{icomma}
\usepackage{xcolor}
\usepackage{tikz}
\usepackage{pgfplots}
\usepackage[most]{tcolorbox}
\usepackage{enumitem}
\usepackage[a4paper,margin=2cm,headheight=26pt,headsep=10pt]{geometry}
\usepackage{fancyhdr}
\usepackage[hidelinks]{hyperref}
\usetikzlibrary{arrows.meta,calc}
\pgfplotsset{compat=1.18}
\definecolor{primary}{RGB}{222,49,99}
\definecolor{primarydark}{RGB}{150,22,55}
\definecolor{methcol}{RGB}{52,92,148}
\definecolor{excol}{RGB}{0,135,90}
\definecolor{defcol}{RGB}{0,114,178}
\definecolor{curveA}{RGB}{0,90,200}
\definecolor{curveB}{RGB}{210,30,40}
\tcbset{boxbase/.style={enhanced,breakable,boxrule=0.9pt,arc=2.5pt,
  left=3mm,right=3mm,top=2.2mm,bottom=2.2mm,fonttitle=\bfseries,coltitle=white}}
\newtcolorbox{cle}[1][]{boxbase,colback=defcol!6,colframe=defcol,
  title={\textsc{À retenir}\ifx\relax#1\relax\relax\else\textmd{ : \itshape #1}\fi}}
\newtcolorbox{methode}[1][]{boxbase,colback=methcol!7,colframe=methcol,sharp corners,
  title={\textsc{Méthode}\ifx\relax#1\relax\relax\else\textmd{ --- \itshape #1}\fi}}
\newtcolorbox{exemplebox}[1][]{boxbase,colback=excol!5,colframe=excol,
  title={\textsc{Exemple}\ifx\relax#1\relax\relax\else\textmd{ : \itshape #1}\fi}}
\newcommand{\R}{\mathbb{R}}
\newcommand{\ds}{\displaystyle}
\pagestyle{fancy}\fancyhf{}
\fancyhead[L]{\small\textcolor{primarydark}{\textbf{Thème 4} — Systèmes d'équations}}
\fancyhead[R]{\small\textcolor{primarydark}{\textsc{Tutoriel}}}
\fancyfoot[C]{\small\thepage}
\renewcommand{\headrulewidth}{0.8pt}
\begin{document}

\begin{center}
{\LARGE\bfseries\color{primarydark} Thème 4 — Systèmes d'équations}\\[2pt]
{\color{primary}\rule{0.5\textwidth}{1.2pt}}\\[4pt]
{\small\itshape Tutoriel — la mécanique à maîtriser avant les exercices}
\end{center}

\vspace{1mm}
\noindent Un \textbf{système} regroupe plusieurs équations à satisfaire \emph{simultanément}. Sa
solution, c'est le jeu de valeurs qui convient à \emph{toutes} les équations à la fois.
Géométriquement (2 inconnues), chaque équation est une droite et la solution est leur
\textbf{point d'intersection}.

\begin{minipage}{0.56\textwidth}
\begin{cle}[deux méthodes à 2 inconnues]
\textbf{Substitution} : isoler une inconnue dans une équation, la remplacer dans l'autre.\\[2pt]
\textbf{Combinaison} : multiplier les équations pour rendre \emph{opposés} les coefficients d'une
inconnue, puis additionner pour l'éliminer.
\end{cle}
\end{minipage}\hfill
\begin{minipage}{0.4\textwidth}
\centering
\begin{tikzpicture}
\begin{axis}[width=5.6cm,height=4.4cm,axis lines=middle,
 xmin=-0.5,xmax=8.5,ymin=-2.5,ymax=6,xtick=\empty,ytick=\empty,
 xlabel={\scriptsize$x$},ylabel={\scriptsize$y$}]
 \addplot[curveA,line width=1pt,domain=0:8.3]{(5-x)/2};
 \addplot[curveB,line width=1pt,domain=3.8:8]{20-3*x};
 \fill (axis cs:7,-1) circle (2pt);
 \node[font=\scriptsize,anchor=south west] at (axis cs:6.2,-1){$(7,-1)$};
\end{axis}
\end{tikzpicture}\\[-1mm]
{\scriptsize Intersection de $x+2y=5$ et $3x+y=20$.}
\end{minipage}

\begin{exemplebox}[le même système, par les deux méthodes]
$\begin{cases} x+2y=5\\ 3x+y=20\end{cases}$\\[2pt]
\textbf{Substitution :} $x=5-2y$, puis $3(5-2y)+y=20\Rightarrow 15-5y=20\Rightarrow y=-1$, d'où $x=7$.\\
\textbf{Combinaison :} $3\times$(1)$-$(2) : $\ (3x+6y)-(3x+y)=15-20\Rightarrow 5y=-5\Rightarrow y=-1$, puis $x=7$.\\
Dans les deux cas $S=\{(7\,;-1)\}$.
\end{exemplebox}

\begin{methode}[3 inconnues : le pivot de Gauss]
On rend le système « triangulaire » en éliminant les inconnues une à une, par des opérations qui ne
changent pas les solutions ($L_i\leftarrow L_i-k\,L_j$) :
\begin{enumerate}[leftmargin=1.6em,topsep=1pt,itemsep=0.5pt]
  \item garder $L_1$, s'en servir pour éliminer $x$ dans $L_2$ et $L_3$ ;
  \item garder la nouvelle $L_2$, éliminer une 2\textsuperscript{e} inconnue dans $L_3$ ;
  \item $L_3$ ne contient plus qu'une inconnue : on la calcule, puis on \emph{remonte}.
\end{enumerate}
\end{methode}

\begin{exemplebox}[pivot de Gauss $3\times3$]
$\begin{cases} x+2y+z=4\ (L_1)\\ 2x+y-z=5\ (L_2)\\ 3x+y+2z=3\ (L_3)\end{cases}$
$\xrightarrow[L_3\leftarrow L_3-3L_1]{L_2\leftarrow L_2-2L_1}$
$\begin{cases} x+2y+z=4\\ -3y-3z=-3\\ -5y-z=-9\end{cases}$\\[3pt]
Puis $L_3\leftarrow L_2-3L_3$ donne $12y=24\Rightarrow y=2$ ; on remonte : $z=-1$ puis $x=1$.
Donc $S=\{(1\,;2\,;-1)\}$.
\end{exemplebox}

\begin{methode}[traduire un énoncé en système]
(1) \textbf{Nommer} les inconnues ; (2) \textbf{traduire} chaque phrase en une équation (« le triple de »,
« augmenté de », « la somme »\dots) ; (3) \textbf{nettoyer} (chasser les fractions en multipliant) ;
(4) \textbf{résoudre} et \textbf{interpréter} le résultat dans le contexte.
\end{methode}

\end{document}
